% =====================================================================
% DOCUMENTCLASS
% =====================================================================
\documentclass[12pt,oneside]{book}

% =====================================================================
% METADATOS
% =====================================================================
\newcommand{\universidad}{Universidad de Buenos Aires}
\newcommand{\facultad}{Facultad de Derecho}
\newcommand{\materia}{Derecho Civil}
\newcommand{\titulo}{Introducción a los Actos Jurídicos}
\newcommand{\autor}{Isaac Edgar Camacho Ocampo}
\newcommand{\anio}{2026}

% =====================================================================
% PAQUETES
% =====================================================================
\usepackage[utf8]{inputenc}
\usepackage[T1]{fontenc}
\usepackage[spanish,es-nodecimaldot]{babel}

\usepackage[
    top=2.5cm,
    bottom=2.5cm,
    left=3cm,
    right=2.5cm,
    headheight=15pt
]{geometry}

\usepackage{amsmath}
\usepackage{amssymb}
\usepackage{mathtools}

\usepackage{graphicx}
\usepackage{float}
\usepackage{caption}
\usepackage{subcaption}

\usepackage{booktabs}
\usepackage{array}
\usepackage{longtable}
\usepackage{tabularx}

\usepackage{enumitem}

\usepackage{xcolor}
\usepackage[most]{tcolorbox}

\usepackage{fancyhdr}
\usepackage{ragged2e}

\usepackage[
    colorlinks=true,
    linkcolor=blue!50!black,
    citecolor=green!40!black,
    urlcolor=blue!60!black,
    pdftitle={\titulo},
    pdfauthor={\autor},
    pdfsubject={Apuntes universitarios},
    bookmarks=true
]{hyperref}

% =====================================================================
% CONFIGURACIÓN
% =====================================================================
\graphicspath{
    {./img/}
    {./graficos/}
    {./figuras/}
}

\setcounter{tocdepth}{2}

\setlength{\parindent}{0pt}
\setlength{\parskip}{0.5em}

\setlist[itemize]{
    topsep=0.3em,
    itemsep=0.2em
}

\setlist[enumerate]{
    topsep=0.3em,
    itemsep=0.2em
}

\clubpenalty=10000
\widowpenalty=10000

% =====================================================================
% COMANDOS
% =====================================================================
\newcommand{\abs}[1]{\left|#1\right|}
\newcommand{\norm}[1]{\left\lVert#1\right\rVert}
\newcommand{\vect}[1]{\mathbf{#1}}
\newcommand{\deriv}[2]{\frac{d#1}{d#2}}
\newcommand{\pderiv}[2]{\frac{\partial#1}{\partial#2}}

% =====================================================================
% ENTORNOS / CAJAS
% =====================================================================
\newtcolorbox{definition}[1]{
    enhanced,
    breakable,
    title={Definición: #1},
    fonttitle=\bfseries,
    colback=blue!3,
    colframe=blue!50!black,
    boxrule=0.8pt,
    arc=2mm
}

\newtcolorbox{important}{
    enhanced,
    breakable,
    title={Importante},
    fonttitle=\bfseries,
    colback=yellow!8,
    colframe=orange!70!black,
    boxrule=0.8pt,
    arc=2mm
}

\newtcolorbox{example}[1]{
    enhanced,
    breakable,
    title={Ejemplo: #1},
    fonttitle=\bfseries,
    colback=green!4,
    colframe=green!40!black,
    boxrule=0.8pt,
    arc=2mm
}

\newtcolorbox{formula}[1]{
    enhanced,
    breakable,
    title={Fórmula / Regla: #1},
    fonttitle=\bfseries,
    colback=purple!4,
    colframe=purple!50!black,
    boxrule=0.8pt,
    arc=2mm
}

\newtcolorbox{exercise}[1]{
    enhanced,
    breakable,
    title={Ejercicio: #1},
    fonttitle=\bfseries,
    colback=gray!5,
    colframe=gray!60!black,
    boxrule=0.8pt,
    arc=2mm
}

\newtcolorbox{note}{
    enhanced,
    breakable,
    title={Nota},
    fonttitle=\bfseries,
    colback=white,
    colframe=gray!50,
    boxrule=0.6pt,
    arc=2mm
}

% =====================================================================
% CONFIGURACIÓN FINAL (encabezados y pies)
% =====================================================================
\pagestyle{fancy}
\fancyhf{}

\fancyhead[L]{\nouppercase{\leftmark}}
\fancyhead[R]{\materia}

\fancyfoot[C]{\thepage}

\renewcommand{\headrulewidth}{0.4pt}
\renewcommand{\footrulewidth}{0pt}

\fancypagestyle{plain}{
    \fancyhf{}
    \fancyfoot[C]{\thepage}
    \renewcommand{\headrulewidth}{0pt}
}

% =====================================================================
% DOCUMENT
% =====================================================================
\begin{document}

% ---------------------------------------------------------------------
% PORTADA
% ---------------------------------------------------------------------
\begin{titlepage}

\thispagestyle{empty}

\begin{center}

\vspace*{1cm}

{\large \textsc{\universidad}}

\vspace{0.2cm}

{\large \textsc{\facultad}}

\vspace{3cm}

{\Huge \textbf{\materia}}

\vspace{1cm}

{\LARGE \titulo}

\vspace{0.8cm}

\textit{Comprender los elementos del acto jurídico es el primer paso para transformar la teoría civil en criterio profesional.}

\vspace{1.2cm}

\rule{0.70\textwidth}{0.5pt}

\vspace{0.5cm}

{\large Apuntes de estudio}

\vspace{0.5cm}

\rule{0.70\textwidth}{0.5pt}

\vfill

{\large \autor}

\vspace{0.3cm}

{\large \anio}

\end{center}

\vfill

\begin{flushleft}
\textit{Este es un modesto aporte a los estudiantes de ``\materia'' no pretende sustituir el material oficial de la materia.}
\end{flushleft}

\end{titlepage}

% ---------------------------------------------------------------------
% ÍNDICE
% ---------------------------------------------------------------------
\tableofcontents

% =====================================================================
% CONTENIDO
% =====================================================================

\chapter{Fundamentos y definición del acto jurídico}

\section{Clasificación de los hechos jurídicos}

El estudio del acto jurídico se inscribe dentro de la teoría general de los hechos jurídicos. Retomando un cuadro de clasificación presentado con anterioridad, los hechos jurídicos se dividen del siguiente modo:

\begin{itemize}
    \item Hechos jurídicos de la naturaleza y hechos humanos.
    \item Dentro de los hechos humanos: \textbf{voluntarios} e \textbf{involuntarios}. Los voluntarios son aquellos ejecutados con discernimiento, intención y libertad; los involuntarios son los que carecen de discernimiento, intención o libertad.
    \item Dentro de los hechos humanos voluntarios: \textbf{actos lícitos} e \textbf{ilícitos} (estos últimos comprenden los delitos y los cuasidelitos).
    \item Dentro de los actos lícitos: \textbf{simples actos lícitos} y \textbf{actos jurídicos}.
\end{itemize}

Los actos jurídicos constituyen, dentro de esta clasificación, la categoría central del presente estudio: son hechos humanos voluntarios lícitos que tienen por fin inmediato producir la adquisición, modificación o extinción de derechos, relaciones o situaciones jurídicas.

\section{Ubicación normativa del acto jurídico}

Todo lo referido a los hechos y actos jurídicos se encuentra ubicado en el \textbf{Libro Primero, Título Cuarto} del Código Civil y Comercial. Esta ubicación representa un avance metodológico significativo respecto del Código de Vélez, que regulaba estos temas en la Sección Segunda del Libro Segundo (el libro de los derechos personales). El nuevo Código traslada estos contenidos a la denominada \textbf{parte general}, reconociendo así su carácter estructural para toda la teoría del derecho civil.

\section{Acto jurídico y contrato: alcance de la teoría general}

El Código Civil y Comercial argentino contiene un sistema general de los actos jurídicos, a diferencia del Código francés, que no trata los actos jurídicos en general sino que se limita al contrato. \textbf{El acto jurídico es un concepto más amplio que el contrato}; el contrato es una especie dentro del género acto jurídico.

\begin{definition}{Contrato}
El contrato es el acto jurídico en el cual hay acuerdo de voluntades para el surgimiento de obligaciones en el campo patrimonial.
\end{definition}

\begin{example}{Contratos como especie de acto jurídico}
\begin{itemize}
    \item \textbf{Contrato de compraventa}: acuerdo de voluntades para transferir la propiedad de una cosa contra el pago de un precio.
    \item \textbf{Contrato de locación}: acuerdo para transferir la tenencia de un bien —mueble o inmueble— contra la entrega de una renta mensual (canon).
\end{itemize}
\end{example}

No todos los actos jurídicos son contratos. Existen actos jurídicos que no responden a esta estructura de acuerdo patrimonial de voluntades:

\begin{itemize}
    \item \textbf{Testamento}: acto que tiene por finalidad disponer de los bienes para el supuesto de la propia muerte, asignándolos según la voluntad del testador. Produce efectos luego de la muerte, pero no constituye un contrato.
    \item \textbf{Matrimonio}: acto jurídico de carácter especial que crea relaciones jurídicas de naturaleza personal.
\end{itemize}

\section{Definición y características del acto jurídico}

\begin{definition}{Acto jurídico (Artículo 259, Código Civil y Comercial de la Nación)}
Acto jurídico es el acto voluntario lícito que tiene por fin inmediato la adquisición, modificación o extinción de derechos, relaciones o situaciones jurídicas.
\end{definition}

De esta definición se desprenden tres elementos fundamentales: que el acto sea \textbf{voluntario}, que sea \textbf{lícito}, y que tenga por \textbf{fin inmediato} producir consecuencias jurídicas.

\begin{important}
La idea de ``fin inmediato'' es crucial: significa que el acto jurídico busca de manera directa la producción de consecuencias jurídicas. Esto es lo que lo diferencia del simple acto lícito, que puede producir consecuencias jurídicas, pero sin que ello constituya su finalidad.
\end{important}

\section{Acto jurídico frente al simple acto lícito}

La distinción entre el acto jurídico y el simple acto lícito reside en la intención del sujeto respecto del efecto jurídico producido.

\begin{example}{El pescador}
Un pescador que sale a pescar, al apropiarse de lo pescado, realiza un simple acto lícito. Se produce un efecto jurídico —el pescador pasa a ser dueño de lo que pescó—, pero su intención original al pescar no es adquirir propiedad, sino pescar. El efecto jurídico es una consecuencia accesoria o secundaria del acto, no su fin inmediato.

Por el contrario, cuando se celebra un contrato de compraventa, el fin inmediato del sujeto es adquirir la propiedad de una cosa. El efecto jurídico es precisamente el objetivo buscado, la intención directa. Por ello, se trata de un acto jurídico.
\end{example}

\chapter{Elementos del acto jurídico}

El acto jurídico se compone de los siguientes elementos esenciales:

\begin{enumerate}
    \item \textbf{Sujeto}: la persona que concurre al acto para su celebración.
    \item \textbf{Objeto}: la materia sobre la que recae el acto, es decir, los bienes o hechos que constituyen su contenido.
    \item \textbf{Forma}: la manera en que se exterioriza el acto.
    \item \textbf{Causa}: la finalidad perseguida por las partes al realizar el acto.
\end{enumerate}

% TODO: contenido incompleto o ambiguo en las notas originales — el tratamiento detallado del elemento "sujeto" y del elemento "forma" no fue desarrollado en esta clase; el apunte indica expresamente que la forma se abordará en una clase posterior, y el sujeto no vuelve a ser tratado en el material disponible.

El presente material se concentra en el desarrollo del \textbf{objeto} y de la \textbf{causa} del acto jurídico. La forma será abordada en una clase posterior. Con carácter ulterior corresponderá también el estudio de los elementos accidentales o modalidades del acto jurídico (condición, plazo y cargo), que pueden estar o no presentes en un acto jurídico determinado.

\chapter{El objeto del acto jurídico}

\section{Fundamento normativo}

El artículo 279 del Código Civil y Comercial es central en materia de objeto de los actos jurídicos. Este artículo reconoce como antecedente el artículo 953 del Código Civil de Vélez, disposición que tuvo enorme influencia a lo largo de toda la vigencia del código anterior, en tanto marcaba el contenido moral del acto.

\begin{definition}{Objeto del acto jurídico (Artículo 279, Código Civil y Comercial de la Nación)}
El objeto del acto jurídico debe ser un hecho o un bien. Los hechos no pueden ser imposibles (ni jurídica ni prácticamente), prohibidos por la ley, o contrarios a la moral y a las buenas costumbres. Los bienes que son objeto del acto no pueden ser prohibidos, ni ser de carácter extrapatrimonial a menos que exista una autorización legal.
\end{definition}

\section{Contenido moral y límite de la autonomía de la voluntad}

Un acto jurídico debe ser conforme a la ley, a la moral y al orden público. No existe una libertad absoluta en la autonomía de la voluntad: esta se encuentra limitada y moldeada por razones de distinto tipo —la ley, el orden público, la moral y la dignidad de la persona.

\begin{note}
A lo largo de la historia, la disposición del artículo 953 del Código de Vélez originó numerosas limitaciones y aplicaciones jurisprudenciales, entre ellas:
\begin{itemize}
    \item Antes de la sanción de la Ley 17.711, dicha disposición dio lugar a la aplicación del vicio de \textbf{lesión}, que permitía anular un acto cuando existía una notable desproporción entre las prestaciones.
    \item Limitación en la aplicación de \textbf{cláusulas penales}, entendidas como la cláusula que se impone a quien incumple un contrato, obligándolo a pagar, por ejemplo, intereses punitorios.
    \item Limitación de los intereses cuando resultan \textbf{usurarios}: cobrar una tasa de interés excesiva o desproporcionada constituye usura, considerada inmoral. Aunque ello depende de la situación y el contexto (por ejemplo, en escenarios inflacionarios), cobrar tasas del 200 al 300 por ciento de interés es claramente inmoral.
\end{itemize}
\end{note}

Existe una tensión ideológica sobre este punto: una mirada más liberal sostiene que cada parte es libre de pactar las tasas de interés que desee. Sin embargo, en Argentina —y en la mayoría de las legislaciones modernas— aun cuando no exista un límite normativo expreso a los intereses, el objeto del acto jurídico debe ser conforme a la moral y a las buenas costumbres, principio que conserva plena vigencia en la actualidad.

\section{Requisitos del objeto}

El artículo 279 divide el objeto del acto jurídico en dos grandes categorías: los hechos y los bienes.

\subsection{Los hechos como objeto del acto}

Los hechos, en tanto acontecimientos, no pueden ser:

\begin{enumerate}
    \item \textbf{Imposibles} (ni jurídica ni prácticamente): no puede pactarse lo imposible. Cruzar el océano Atlántico nadando, por ejemplo, no puede ser objeto de un acto jurídico por tratarse de un hecho prácticamente imposible.
    \item \textbf{Prohibidos por la ley}: si la ley prohíbe determinado hecho, este no puede ser objeto de un acto jurídico.
    \item \textbf{Contrarios a la moral y a las buenas costumbres}: constituye una exigencia fundamental que el acto jurídico se adecue a los principios morales básicos.
\end{enumerate}

\subsection{Los bienes como objeto del acto}

El artículo 279 hace referencia genérica a prohibiciones que pueden surgir de la ley. Así, por ejemplo, ciertos bienes no pueden ser objeto de determinados contratos: los bienes de dominio público no pueden ser objeto de un acto jurídico de disposición.

\section{Objeto ilícito: el caso de la venta de influencias}

\begin{example}{Venta de influencias y acto de corrupción}
Una persona se compromete a realizar tratativas para obtener un privilegio en sede administrativa. En el fondo, se trata de una venta de influencias, un acto de corrupción.

Este acto tiene un objeto ilícito porque vender influencias es contrario a la ética pública y constituye un acto de corrupción que afecta el orden público.

Si posteriormente esa persona pretende hacer valer la obligación —por ejemplo, cobrar lo que le fue prometido—, la consecuencia es la \textbf{nulidad del acto}. La nulidad surge precisamente porque el objeto no responde a los requisitos del artículo 279. Se trata de una nulidad de tipo \textbf{absoluta}, regulada en el artículo 386 del Código Civil y Comercial.
\end{example}

\chapter{La causa del acto jurídico}

\section{Causa fuente y causa fin}

El concepto de causa admite varias acepciones en el derecho civil. En particular, se distingue entre la causa fuente y la causa fin.

\begin{definition}{Causa fuente}
Es la fuente de las obligaciones, es decir, el origen del cual surgen. Las obligaciones surgen de los contratos, en general de los actos jurídicos, y también de los actos ilícitos. No puede haber una obligación sin una causa, sin un elemento que le dé origen.
\end{definition}

\begin{definition}{Causa fin}
Es el fin perseguido por las partes al momento de celebrar un acto: qué se propusieron las partes al celebrarlo y cuál es la finalidad que persiguen.
\end{definition}

El desarrollo que sigue se concentra en el problema de la \textbf{causa fin}, y no en la causa fuente.

\section{Evolución histórica de la teoría de la causa}

El concepto de causa ha sido objeto de debate a lo largo de la historia del derecho, dando lugar a distintas posturas doctrinarias. Elementos vinculados a la valoración de los actos aparecen desde la antigüedad, especialmente con los canonistas y los glosadores.

\subsection{Teoría causalista}

Esta teoría, atribuida a un autor francés, sostiene que la causa fin de un contrato es la prestación de la otra parte. Se presenta como algo evidente que, al celebrar un contrato sinalagmático —perfectamente bilateral, en el que una parte tiene una prestación que la obliga y la otra, a su vez, otra prestación que también la obliga—, cada parte persigue como causa la prestación de la contraria.

\begin{example}{Compraventa bajo la teoría causalista}
En una compraventa, la causa por la cual una parte celebra el contrato es adquirir la cosa, y para la otra parte, la causa es recibir el precio.
\end{example}

\subsection{Crítica anticausalista}

Un autor belga —mencionado posteriormente por otros autores, entre ellos Planiol— critica la tesis causalista. El argumento principal sostiene que, si la causa fin es la prestación de la otra parte, entonces la causa se identifica con el objeto mismo del contrato, sin ofrecer ninguna novedad explicativa.

La crítica anticausalista señala que la idea de causa fin no agrega nada que no esté ya detrás del objeto: si el objeto es la transferencia de la propiedad (en una compraventa) y la causa es adquirir esa propiedad, causa y objeto se superponen. La postura anticausalista afirma, en consecuencia, que la causa no constituye un elemento independiente.

\subsection{Teoría moderna de la causa}

Enrique Capitant, hacia comienzos del siglo XX, retoma la idea de la causa y la profundiza. Para Capitant, la causa relevante es la finalidad perseguida por las partes, una finalidad que estas incorporan y que buscan en el momento de celebrar el acto.

\section{La causa en el Código de Vélez}

El Código Civil de Vélez no menciona la causa entre los elementos del acto de manera explícita en la sección específica de actos jurídicos. Sin embargo, el tema de la causa aparecía en algunos artículos ubicados en la parte de obligaciones (artículos 499 a 502), en el Libro Segundo, que se abría con la sección de obligaciones.

\begin{definition}{Artículo 499, Código Civil de Vélez}
No hay obligación sin causa. Se entiende por tal el hecho o acto lícito o ilícito en la relación de familia, o entre extraños, que es la causa generadora del derecho que faculta al acreedor a exigir del deudor la prestación debida.
\end{definition}

Este artículo se refiere claramente a la \textbf{causa fuente} (el hecho o acto que genera la obligación). Sin embargo, en los artículos 500 y 502, y especialmente en las notas que Vélez Sarsfield colocaba al pie de la ley —donde se critica la confusión entre causa fuente y causa fin—, aparece la idea de causa fin vinculada a la finalidad perseguida.

\begin{definition}{Artículo 502, Código Civil de Vélez (principio de la causa ilícita)}
De ningún efecto es el acto constituido con una causa ilícita, si esa causa es contraria a las leyes o al orden público.
\end{definition}

Aunque la causa no estaba expresamente mencionada como elemento en el Código de Vélez, existía la presunción de que, si las partes no expresaban la causa, esta debía existir igualmente, y la obligación se mantenía válida aun cuando la causa expresada fuera falsa.

\begin{important}
Una obligación con causa ilícita carece de todo efecto. Si la causa es ilícita, si contraviene las leyes o el orden público, el acto es nulo.
\end{important}

% TODO: contenido incompleto o ambiguo en las notas originales — el apunte menciona una observación del docente sobre el "artículo 429 del Código de Vélez", indicando que, cuando dicho artículo se refiere a una causa ilícita, en realidad remite a la causa fuente y no a la causa fin (ejemplo dado: quien causa un daño a otro debe repararlo, lo cual es causa fuente de la obligación de reparar). El apunte no precisa con exactitud el contenido ni la ubicación de dicho artículo dentro del sistema de Vélez, por lo que se conserva la referencia tal como fue consignada en clase.

\section{La causa en el Código Civil y Comercial}

El nuevo Código Civil y Comercial incorpora expresamente la causa como elemento del acto jurídico. El artículo 281 es la disposición central de esta regulación.

\begin{definition}{Artículo 281, Código Civil y Comercial}
La causa es el fin inmediato perseguido por las partes que ha sido determinante de la voluntad. Cuando las partes no la expresan, se presume que existe. Se presume también su licitud, a menos que resulte ilícita del acto mismo. También integran la causa los motivos exteriorizados cuando sean lícitos y hayan sido incorporados al acto en forma expresa o tácita, y sean esenciales para ambas partes.
\end{definition}

En rigor, lo que resulta relevante de esta definición es que la causa es el fin perseguido por las partes que ha sido determinante de la voluntad.

\subsection{Dimensiones objetiva y subjetiva de la causa}

El artículo 281 incorpora dos dimensiones de la causa:

\begin{enumerate}
    \item \textbf{Dimensión objetiva}: el fin inmediato autorizado por el ordenamiento jurídico. Es la finalidad típica que el derecho reconoce como válida para ese tipo de acto. En una compraventa, la causa objetiva es transferir la propiedad de una cosa a cambio de un precio.
    \item \textbf{Dimensión subjetiva}: los motivos o intenciones de las partes. Estos integran la causa cuando son exteriorizados (expresados o manifestados), lícitos, y esenciales para ambas partes; es decir, cuando ambas partes conocen y aceptan esos motivos como parte de la razón por la cual celebran el acto.
\end{enumerate}

\begin{example}{El caso de la profesora de francés}
Una profesora de francés contrata a un escritor para que redacte un ensayo que ella presentará, traducido, en un concurso destinado a obtener una beca para viajar a Francia. La profesora presentará el ensayo como propio.

En este caso se distingue con claridad entre objeto y causa:
\begin{itemize}
    \item El \textbf{objeto} del acto es contratar a alguien para que realice un ensayo o un trabajo. Esto es, en sí mismo, perfectamente válido y lícito.
    \item La \textbf{causa} (finalidad perseguida) es presentar el ensayo como propio para engañar a un concurso y obtener una beca. Esto es ilícito.
\end{itemize}

Si una parte persigue una causa ilícita pero la otra parte lo ignora, esta última no resulta afectada: el motivo ilícito no integra el acto. Pero si se interpreta que ambas partes conocen la causa y esta es esencial para ambas, entonces dicha causa integra el acto y lo vicia.
\end{example}

\subsection{Síntesis de las características de la causa}

\begin{enumerate}
    \item El Código Civil y Comercial incorpora la causa como elemento del acto jurídico en el artículo 281.
    \item La causa comprende dos dimensiones: una \textbf{objetiva} (el fin inmediato autorizado por el ordenamiento, determinante de la voluntad) y una \textbf{subjetiva} (los motivos o intenciones de las partes).
    \item Solo integran la causa los motivos o intenciones cuando: (a) son exteriorizados, (b) son lícitos, y (c) son esenciales para ambas partes.
    \item Quedan excluidos de la causa los puros motivos o intenciones personales que no estén exteriorizados ni sean esenciales para ambas partes.
\end{enumerate}

\section{Régimen de la causa}

\begin{definition}{Artículo 282, Código Civil y Comercial}
Cuando la causa no está expresada, se presume que existe. Se presume también su licitud, a menos que resulte ilícita del acto mismo.
\end{definition}

El artículo 282 retoma los principios ya contenidos en los artículos 500, 501 y 502 del Código de Vélez:

\begin{itemize}
    \item Si la causa no está expresada, se presume que existe (presunción de causa).
    \item El acto vale aunque la causa expresada sea falsa. Si lo que se aparenta es falso pero existe una causa verdadera, se estará ante una simulación, aunque el acto siga siendo válido.
    \item Una causa ilícita vicia el acto: si la causa es ilícita —contraria a las leyes y al orden público—, el acto es nulo o de ningún efecto.
\end{itemize}

\section{Actos abstractos}

\begin{definition}{Artículo 1014, Código Civil y Comercial (en materia de contratos)}
La inexistencia, falsedad o ilicitud de la causa no son discutibles en el acto abstracto.
\end{definition}

% TODO: contenido incompleto o ambiguo en las notas originales — el apunte atribuye esta regla al "artículo 283" en el desarrollo oral de la clase, mientras que el recuadro normativo transcripto en el propio material la identifica como "artículo 1014". Se conserva la referencia tal como figura en las notas originales, sin resolver la discrepancia mediante conocimiento externo.

Un \textbf{acto abstracto} es aquel que se independiza de su causa y recibe un tratamiento autónomo. El caso típico es el pagaré.

\begin{example}{El pagaré}
Una persona firma un pagaré porque contrae una deuda con un hospital. Al ingresar al hospital, el pagaré adquiere autonomía propia: puede circular y transferirse a terceros. Si posteriormente esa persona ya pagó al hospital y alguien pretende ejecutar el pagaré, no puede discutirse la causa (el hecho de que ya se pagó). La causa por la cual se emitió el pagaré no se discute en los títulos abstractos.
\end{example}

\section{Cierre de la clase y próximos temas}

En esta clase se desarrollaron los elementos del acto jurídico, con especial énfasis en el objeto y en la causa. En la clase siguiente se profundizará en la clasificación del objeto y, probablemente, se avanzará sobre las modalidades de los actos jurídicos (condición, plazo y cargo).

% =====================================================================
% CORNELL
% =====================================================================
\chapter{Repaso mediante el método Cornell}

El siguiente cuadro reproduce, en formato de tabla Cornell, las preguntas y respuestas centrales para el repaso activo de los contenidos desarrollados en este material.

\begin{center}
\begin{longtable}{
    >{\raggedright\arraybackslash}p{0.28\textwidth}
    >{\raggedright\arraybackslash}p{0.62\textwidth}
}
\toprule
\textbf{Preguntas / palabras clave} & \textbf{Notas / respuestas} \\
\midrule
\endfirsthead
\toprule
\textbf{Preguntas / palabras clave} & \textbf{Notas / respuestas} \\
\midrule
\endhead
\midrule
\multicolumn{2}{r}{\textit{Continúa en la página siguiente}} \\
\midrule
\endfoot
\bottomrule
\endlastfoot

\textbf{¿Qué es un acto jurídico?} &
Es el acto voluntario lícito con fin inmediato de producir la adquisición, modificación o extinción de derechos, relaciones o situaciones jurídicas (art.\ 259, Código Civil y Comercial). Lo fundamental es que la intención directa de generar consecuencias jurídicas distingue al acto jurídico del simple acto lícito. \\

\addlinespace

\textbf{¿Cuál es la diferencia entre acto jurídico y simple acto lícito?} &
En el acto jurídico, la producción de consecuencias jurídicas es el fin inmediato e intencional. En el simple acto lícito, las consecuencias jurídicas se producen, pero no son el fin buscado (por ejemplo, el pescador que se apropia de lo que pesca). Lo determinante es la intención o finalidad del sujeto respecto del efecto jurídico. \\

\addlinespace

\textbf{¿Cuáles son los elementos del acto jurídico?} &
1.\ Sujeto: la persona que realiza el acto. 2.\ Objeto: la materia o los bienes sobre los que recae el acto. 3.\ Forma: la manera de exteriorizar el acto. 4.\ Causa: la finalidad perseguida por las partes. \\

\addlinespace

\textbf{¿Qué requisitos debe cumplir el objeto?} &
Debe ser posible (física y jurídicamente); no debe estar prohibido por la ley; debe ser lícito, es decir, conforme a la moral, a las buenas costumbres y al orden público; y no debe ser de carácter extrapatrimonial, salvo autorización legal. \\

\addlinespace

\textbf{¿Qué es la causa del acto?} &
Es el fin inmediato perseguido por las partes que ha sido determinante de la voluntad (art.\ 281, Código Civil y Comercial). Incorpora una dimensión objetiva (el fin autorizado por el derecho) y una dimensión subjetiva (los motivos lícitos, exteriorizados y esenciales para ambas partes). Se distingue de la causa fuente (origen de la obligación), que resulta menos relevante a estos efectos. \\

\addlinespace

\textbf{¿Qué diferencia hay entre objeto ilícito y causa ilícita?} &
En el objeto ilícito, la materia misma sobre la que recae el acto es ilícita (por ejemplo, vender un bien prohibido o contratar un acto ilegal). En la causa ilícita, el objeto es lícito, pero la finalidad o el motivo que persiguen las partes es ilícito (por ejemplo, contratar la redacción de un ensayo cuyo fin es el fraude académico). Ambas vician el acto, aunque por razones distintas. \\

\addlinespace

\textbf{¿Qué es un acto abstracto?} &
Es el acto que se independiza de su causa y recibe un tratamiento autónomo. El ejemplo típico es el pagaré: una vez emitido, puede circular sin vincularse a la causa original, y dicha causa no puede discutirse al momento de su ejecución. La inexistencia, falsedad o ilicitud de la causa no son discutibles en el acto abstracto. \\

\midrule
\multicolumn{2}{p{0.90\textwidth}}{
\textbf{Resumen:} Los actos jurídicos constituyen el fenómeno jurídico fundamental mediante el cual las personas, voluntariamente, crean, modifican o extinguen derechos y obligaciones. Su estudio requiere comprender que un acto jurídico no es simplemente cualquier conducta lícita: su característica esencial es que su fin inmediato es la producción intencional de consecuencias jurídicas. Para que un acto sea válido debe reunir cuatro elementos: un sujeto (persona con capacidad), un objeto (materia lícita, posible y conforme a la moral), una forma (manera de exteriorización) y una causa (finalidad perseguida). El objeto y la causa constituyen los límites más importantes a la autonomía de la voluntad: no puede haber objeto imposible, prohibido o contrario a la moral, ni causa ilícita. La distinción entre objeto ilícito y causa ilícita es crucial, porque permite comprender que un acto puede tener un objeto perfectamente válido pero una finalidad viciada. El Código Civil y Comercial introduce innovaciones relevantes: incorpora la causa expresamente como elemento (art.\ 281), reconoce sus dimensiones objetiva y subjetiva, y establece presunciones de existencia y licitud de la causa. Salvo en los actos abstractos (como el pagaré), la causa es siempre discutible y determinante: una causa ilícita torna nulo el acto, sin importar cuán lícito sea su objeto.
} \\

\end{longtable}
\end{center}

\end{document}
